\documentclass[11pt]{article}

\usepackage[english]{babel}
\usepackage[utf8]{inputenc}
\usepackage{amsmath}
\usepackage{amssymb}
\usepackage{listings}
\usepackage{epsfig}
\usepackage{hyperref}
\usepackage{gb4e}
\usepackage{enumitem}

\lstset{
  frame=single,
  numbers=left,                    % where to put the line-numbers; possible values are (none, left, right)
  numbersep=5pt,                   % how far the line-numbers are from the code
  literate=  {→}{{$\rightarrow$}}1 {↓}{{$\downarrow$}}1 {↑}{{$\uparrow$}}1 {⇓}{{$\Downarrow$}}1 {⇑}{{$\Uparrow$}}1 {⊂}{{$\subset$}}1 {∪}{{$\cup$}}1 {¬}{{$\lnot$}}1 {≠}{{$\neq$}}1 {≤}{{$\leq$}}1 {≥}{{$\geq$}}1 {∨}{{$\lor$}}1 {∧}{{$\land$}}1 {⇒}{{$\Rightarrow$}}1 {⇔}{{$\Leftrightarrow$}}1
  {á}{{\'a}}1 {é}{{\'e}}1 {í}{{\'i}}1 {ó}{{\'o}}1 {ú}{{\'u}}1
  {Á}{{\'A}}1 {É}{{\'E}}1 {Í}{{\'I}}1 {Ó}{{\'O}}1 {Ú}{{\'U}}1
  {à}{{\`a}}1 {è}{{\`e}}1 {ì}{{\`i}}1 {ò}{{\`o}}1 {ù}{{\`u}}1
  {À}{{\`A}}1 {È}{{\'E}}1 {Ì}{{\`I}}1 {Ò}{{\`O}}1 {Ù}{{\`U}}1
  {ä}{{\"a}}1 {ë}{{\"e}}1 {ï}{{\"i}}1 {ö}{{\"o}}1 {ü}{{\"u}}1
  {Ä}{{\"A}}1 {Ë}{{\"E}}1 {Ï}{{\"I}}1 {Ö}{{\"O}}1 {Ü}{{\"U}}1
  {â}{{\^a}}1 {ê}{{\^e}}1 {î}{{\^i}}1 {ô}{{\^o}}1 {û}{{\^u}}1
  {Â}{{\^A}}1 {Ê}{{\^E}}1 {Î}{{\^I}}1 {Ô}{{\^O}}1 {Û}{{\^U}}1
  {Ã}{{\~A}}1 {ã}{{\~a}}1 {Õ}{{\~O}}1 {õ}{{\~o}}1
  {œ}{{\oe}}1 {Œ}{{\OE}}1 {æ}{{\ae}}1 {Æ}{{\AE}}1 {ß}{{\ss}}1
  {ű}{{\H{u}}}1 {Ű}{{\H{U}}}1 {ő}{{\H{o}}}1 {Ő}{{\H{O}}}1
  {ç}{{\c c}}1 {Ç}{{\c C}}1 {ø}{{\o}}1 {å}{{\r a}}1 {Å}{{\r A}}1
  {€}{{\euro}}1 {£}{{\pounds}}1 {«}{{\guillemotleft}}1
  {»}{{\guillemotright}}1 {ñ}{{\~n}}1 {Ñ}{{\~N}}1 {¿}{{?`}}1
}


\title{Elvex -- Documentation}
\author{Lionel Clément}
\date{Version 2.27.0}
%November 2019 - June 2024}

\begin{document}
\maketitle

\section*{Introduction}

\texttt{Elvex} is a surface realisation system based on typed feature
structures, lexical resources and syntagmatic rules. It takes as input a
feature structure representing a concept or a generation request, combines
it with a lexicon and a grammar, and produces one or more textual
realizations.

The central idea of \texttt{Elvex} is to describe generation by means of
rules whose right-hand side symbols exchange information through inherited
and synthesized feature structures. Inherited structures, written
$\downarrow_i$, transmit information from a rule to its components, while
synthesized structures, written $\Downarrow_i$, return information produced
by these components. This makes it possible to express agreement,
selection, optional constituents, alternatives and lexical constraints in a
uniform way.

This document introduces the main elements of the language: feature
structures, lexicons, rules, operational semantics and command-line tools.
It also describes mechanisms used during generation, such as optional
symbols, alternatives, deferred evaluation, tracing and linearization
constraints. The examples are intended to show both the concrete syntax and
the style of grammar writing expected by \texttt{Elvex}.

The documentation is organized progressively. The first sections present
the linguistic and formal notions used by the system. The following
sections describe the syntax of grammars and lexicons, then the operational
statements available inside rules. The final sections explain how to run
\texttt{Elvex}, how to debug grammars, and how to build or consult
compacted lexicons.

\section{Full Example: Getting Started}

An example will familiarize the reader with \texttt{Elvex} before we
go any further. Let us take the case of the text generation of a
nominal phrase in French.

\subsubsection*{Example}

\begin{lstlisting}
NP → det N {
  ↓2 = ↑;
  ↓1 = ⇓2;
  ⇑ = ⇓2;
}

N → adj N {
  [mod:<$Head::$Tail>, $Rest];
  [number:$inhNb] ⊂ ↑;
  ↓1 = [number:$inhNb, gender:$synthGd, $Head];
  ↓2 = [mod:$Tail, $Rest];
  [gender:$synthGd] ⊂ ⇓2;
  ⇑ = ⇓2 ∪ [qual:yes];
}

N → n {
  [mod:NIL];
  ↓1 = ↑;
  ⇑ = ⇓1;
}
\end{lstlisting}

\subsection{Syntagmatic description (\textit{Context-Free Grammar rules})}

The first rule of grammar (line 1) defines \texttt{NP} as
consisting of \texttt{det} and \texttt{N} from left to right.

\texttt{det} is a word; it is defined in the lexicon that we will see
later.

Lines 7) and 16) follow the same logic and we have the
context-free grammar:

\begin{lstlisting}[numbers=none]
NP → det N
N → adj N
N → n
\end{lstlisting}

The only sequence we can generate with this simple example is
a noun, preceded by a determiner and possibly one or several adjectives.

Syntagmatic rules are therefore made to describe the order of
components of phrases, sentences and speeches.

\subsection{Operational semantics}

Each rule contains a set of operations. These operations describe the
functional properties of the language, which may or may not depend on
the order of the words.

\subsubsection*{The basics}

We will explain this semantics through the example given to the
beginning of this section before giving details.

\begin{itemize}
	\item The lines 2) and 3) express the nature of the links that
	      exist between \texttt{NP}, the name \texttt{N} and its determiner
	      \texttt{det}; the main one being that agreement in gender and number
	      occurs between nouns and their specifier.

	      \begin{itemize}
		      \item Line 2) assigns the given nominal phrase properties to the
		            noun. The attribute of the name is noted $\downarrow2$
		            (this is the second term of the right part of the rule), and
		            nominal phrase's is noted $\uparrow$.

		      \item Line 3) retrieves information from analyses of \texttt{N} to
		            assign them to the determiner.  The attribute of the determiner
		            is noted $\downarrow1$ (this is the first term of the right part
		            of the rule), the attribute from the name \texttt{N}
		            is noted $\Downarrow2$.

		      \item Line 4) assigns the analyse of \texttt{N} to the synthesized
		            attribute \texttt{NP}.

	      \end{itemize}

	      With the two first lines, the noun will be chosen in the lexicon
	      according to the concept to be expressed (attribute of the nominal
	      phrase $\uparrow$).  A female singular determiner will be chosen if
	      the gender of the noun is feminine and the nominal phrase is
	      singular.

	      This mechanism, which makes it possible to articulate lexical
	      choices according to the concepts and syntactic properties according
	      to lexical properties is at the heart of \texttt{Elvex} and
	      implements two types of attributes: inherited attributes (noted
	      $\uparrow$, $\downarrow_i$) and synthesized attributes (noted
	      $\Uparrow$, $\Downarrow_i$).

	      Inherited attributes are calculated from left to right of the
	      syntagmatic rules. Synthesized attributes are calculated from right
	      to left. To put it simply, we can say that the inherited attributes
	      convey information from the concept to the lexicon, and that the
	      synthesized attributes convey information from the lexicon to higher
	      layers of the syntactic analysis. Thus inherited attributes allow you to
	      propagate the input to the lexicon to make a lexical choice, whereas
	      conversely, synthesized attributes allow to retrieve the lexical information of
	      the chosen name according to the concept expressed and to propagate
	      this information (grammatical features, selection restrictions,
	      phraseological units constraints, etc.)  to the entire nominal
	      phrase.

	\item The lines 8) to 13) allow you to set the gender and
	      number agreement of the noun and adjective and also to propagate the
	      concepts that produce adjectives. We notice that in French, the
	      number is a semantism expressed directly as a sign associated with
	      the noun, whereas gender is a grammatical feature given by the
	      lexicon. The first will be propagated from the concept to the noun,
	      the second will be propagated from the lexicon to the layers
	      regulating the agreement.


	      \begin{itemize}
		      \item Line 8) is a \textit{guard}, \textit{i. e.} a
		            condition to realize
		            $\mathtt{N} \rightarrow \mathtt{adj}~ \mathtt{N}$. It expresses that the
		            realization of an adjective is conditioned by
		            the existence of a list
		            \verb#<$Head::$Tail># as the value of an attribute \texttt{mod}
		            in the inherited attribute $\uparrow$.

		            \verb#$Head# and
		            \verb#$Tail# are two variables that respectively define the
		            head and the tail of a list, \textit{i. e.} respectively the
		            first item in the list and the rest of the list.

		            if the condition is verified,
		            \verb#$Head#, \verb#$Tail# and
		            \verb#$Rest# will be affected by values given by the
		            subsumption\footnote{A feature structure \texttt{A} subsumes
			            a feature structure \texttt{B} iif all the features of
			            \texttt{B} are present in \texttt{A} with the same values
			            or with values that subsume those of the
			            values of \texttt{A}. } values of \\\verb#[mod:<$Head::$Tail>, $Rest]#
		            by $\uparrow$.

		      \item Line 9) assign to the variable
		            \verb#$inhNb# (\textit{Inherited noun}) the value of the
		            \texttt{number} feature of the inherited attribute
		            $\uparrow$. This allows you to retrieve the number given to the
		            nominal phrase.

		      \item Line 10) assigns to the inherited attribute of the adjective
		            the head of the list \verb#$Head# unified\footnote{The unification
			            of a feature structure \texttt{A} with a feature structure
			            \texttt{B}, is the smallest feature structure \texttt{C} such
			            that \texttt{A} subsumes \texttt{C} and \texttt{B} subsumes
			            \texttt{C}.} with a structure that gives the synthesized gender
		            \verb#$synthGd# (\textit{synthesized gender}) and the inherited
		            number \verb#$inhNb#.

		      \item Line 11) assigns to the inherited attribute of the noun all
		            values passed by inheritance (\verb#$Rest#) and adds the feature
		            \verb#mod# with the value
		            \verb#$Tail#, which is the tail of the list that brings together
		            all other modifiers of the noun (other adjective or a proposal for
		            example).

		      \item Line 12) allows you to assign the variable
		            \verb#$synthGd# which is the gender synthesized from the noun.

		      \item Finally line 13) assigns to the synthesized attribute of the
		            phrase the synthesized attribute of the noun and adds a feature
		            \texttt{qual} which marks that this noun phrase is qualified
		            (which makes it possible to restrict others modifications).

	      \end{itemize}

	\item The rules on lines 17) to 19) are not specific, except for the
	      custody of rule 17) which implements the value \verb#NIL#.

	      \begin{itemize}
		      \item Line 17) is a guard which allows us to prevent the inherited
		            attribute $\uparrow$ to contain any feature
		            \texttt{mod}. \verb#mod:NIL# means that the attribute \verb#mod#
		            does not exist in $\uparrow$.

		      \item Lines 18) and 19) do not require any further comment.

	      \end{itemize}
\end{itemize}

It becomes clear by examining this simple example that the order of
operational rules is not the one given by grammar, but the one
dictated by the availability of operands. Here, a possible order of
application of the rules is 8), 9), 11), 12), 10), and 13).

In the current version of \texttt{Elvex}, this order is calculated from
the availability of operands during rule evaluation.

The advice I can give to develop a grammar is
to write the rules as the synthesized attributes are available.

We see that the operational statements of syntagmatic rules are not
evaluated strictly from left to right or from top to bottom, but according
to the availability of the attributes they depend on.

Before going any further, let us give definitions on the ratings that
allow us to understand the syntax of the rules.

\subsubsection*{Empty rules, optional terms}

It is possible to write empty rules, \textit{i.e.}  components that have no
realization in text production. It is also possible to put optional
elements written between parentheses \verb#()#.

Example: The following rule gives all possible combinations for clitic
pronouns placed before the verb.

\begin{lstlisting}[numbers=none]
  CLITICS → (clr) (cld) (cla) (cld) (clg) (cll) 
\end{lstlisting}

In French, the dative clitic (\texttt{cld}) is sometimes placed
before, sometimes placed after the accusative clitic (\texttt{cla})
according to the person of the pronoun.

\begin{exe}
	\ex \begin{xlist}
		\ex
		\begin{tabbing}
			\textit{Jean} xxx \= \textbf{to\_me}(dative) \= \textbf{it}(accusative) \= \textit{donne}\kill
			\textit{Jean} \>\textit{\textbf{me}} \>\textit{\textbf{le}} \>\textit{donne}\\
			Jean \> \textbf{to\_me}(dative) \> \textbf{it}(accusative) \> gives\\
			Jean gives it to me\\
		\end{tabbing}

		\ex
		\begin{tabbing}
			\textit{Jean} xxx \=\textbf{it}(accusative) \=\textbf{to\_him}(dative) \= \textit{donne}\kill
			\textit{Jean} \>\textit{\textbf{le}} \>\textit{\textbf{lui}} \>\textit{donne}\\
			Jean \>\textbf{it}(accusative) \>\textbf{to\_him}(dative) \>gives\\
			Jean gives it to him.
		\end{tabbing}
	\end{xlist}
\end{exe}

According to this rule, the phrase \texttt{CLITICS} can also be empty,
due to the fact that each term is optional.  A rule whose right-hand side
contains only optional terms may therefore generate the empty string
$\epsilon$. Operationally, an absent optional component is represented by
an empty forest whose span is $[i,i]$: no lexical form is consumed, but the
component is nevertheless considered resolved. This is important because the
rest of the rule may still refer to the presence or absence of this component
with the presence tests described below.

For a rule with $n$ independent optional components, the generator may in
principle produce up to $2^n$ combinations. One of these combinations may be
empty if all optional components are absent. If the output device prints one
realization per line, this empty realization may appear as a blank line or may
not be visually obvious.

\subsubsection*{Paradigms}

Each element of a syntagmatic rule can be a list of terms
separated by the character \verb#|# instead of being a single term. In this
case, the product phrase is constructed with one of the terms in the list.

Example:
The following rule allows you to build a sentence without a subject, with a
nominative clitic subject \texttt{cln} (\textit{je}, \textit{tu},
\textit{il}, $\dots$), or with a subject as a nominal group.

\begin{lstlisting}[numbers=none]
S → (cln | NP) VP
\end{lstlisting}

A grammar written by using massively rules with alternatives and
optional terms is equivalent to a very large grammar.  For example,
the rule describing the order of clitic pronouns is equivalent to a
set of 64 rules:

\begin{lstlisting}[numbers=none]
  CLITICS → 
  CLITICS → clr 
  CLITICS → cld 
  ...
  CLITICS → clr cld 
  CLITICS → clr cld cla
  ... ...
\end{lstlisting}

However, \texttt{Elvex} never builds this set of rules explicitly and
remains efficient with large grammars written with alternatives and
optional terms.

The use of alternatives and optional terms is therefore recommended to
describe the order of words, without regard to the combinatorial
problem.

\paragraph{Operational resolution of optional terms and alternatives}

Optional terms and alternatives are resolved by the generator before the
corresponding component is generated.  An optional term opens two branches:
one branch where the component is absent and one branch where it is present.
When the component is absent, the generator records an empty forest with the
same start and end position. When the component is present, the optional mark
is removed and the component is generated normally.

An alternative such as \verb#NP|cln# also opens one branch for each possible
variant. If the term is both optional and alternative, for example
\verb#(NP|cln)#, the optional decision is made first: the component may be
absent, or it may be present and then one of the alternatives is selected.

These decisions are internal to the generation algorithm, but they are visible
in the operational semantics through the presence tests \verb|#i| and
\verb|#i.j|.

\section{Elvex notations}


\subsection{General syntax}

Elvex source files may contain comments and file inclusions.

\paragraph{Comments}

Two forms of comments are accepted:

\begin{itemize}
	\item line comments, introduced by \verb#//# and ending at the end of the line;
	\item block comments, delimited by \verb#/*# and \verb#*/#.
\end{itemize}

For example:

\begin{lstlisting}[numbers=none]
// This is a line comment

/*
   This is a block comment.
   It may span several lines.
*/
\end{lstlisting}

\paragraph{File inclusion}

An Elvex source file may include another file with the directive
\verb|\#include| at the beginning of a line:

\begin{lstlisting}[numbers=none]
#include path/to/file.elvex
\end{lstlisting}

The included file is read as if its contents occurred at the position of
the directive.  This is useful for splitting large descriptions into several
files.

When verbose mode is enabled, Elvex reports the opening of included files.

A syntagmatic rule is written

$A~ \rightarrow B_1 B_2 \dots B_k <operational~semantics>$

Where $B_i$ is written:

\begin{itemize}
	\item $( B_i )$ if the symbol $B_i$ is optional.
	\item $C_1 | C_2 | \dots | C_n$ if $B_i$ is a choice between the $C_i$
\end{itemize}

$B_i$, $C_i$ can be lexical categories or symbols of the
grammar.

\subsection{Operational Semantics}

\paragraph{Notations}

Let the rule
$A~ \rightarrow B_1 B_2 \dots B_i \dots B_k$

\begin{itemize}
	\item $\uparrow$: refers to the attribute inherited from $A$.
	\item $\downarrow_i$: refers to the attribute inherited from $B_i$.
	\item $\Uparrow$: refers to the synthesized attribute of $A$.
	\item $\Downarrow_i$: refers to the synthesized attribute of $B_i$.
\end{itemize}

\paragraph{ASCII alternatives}

Most symbolic operators also have ASCII alternatives.  These alternatives
are often more convenient when writing grammar files on a standard keyboard.

\begin{center}
\begin{tabular}{lll}
\hline
Symbolic notation & ASCII notation & Meaning \\
\hline
$\rightarrow$ & \verb#-># & rule arrow \\
$\cup$ & \verb#U# & unification \\
$\subset$ & \verb#<<<# & subsumption \\
$\neq$ & \verb#!=# & inequality \\
$\leq$ & \verb#<=# & less than or equal \\
$\geq$ & \verb#>=# & greater than or equal \\
$\lor$ & \verb#||# & logical disjunction \\
$\land$ & \verb#&&# & logical conjunction \\
$\Rightarrow$ & \verb#=># & implication \\
$\Leftrightarrow$ & \verb#<=># & equivalence \\
$\lnot$ & \verb#!# & negation \\
\hline
\end{tabular}
\end{center}

For example, the following two rules use equivalent notations:

\begin{lstlisting}[numbers=none]
S → NP VP { ↓1 = ↑ ∪ [case:nom]; }

S -> NP VP { ↓1 = ↑ U [case:nom]; }
\end{lstlisting}

The named variables are noted by an identifier preceded by ``\$'',
the anonymous variables are noted ``\verb#_#'', the
constant values are noted by identifiers.

The feature structures are noted $[ f_1, f_2, \dots, f_k ]$
where $ f_i$ is a feature as
\begin{itemize}
	\item a linked variable feature
	\item an anonymous variable feature
	\item \verb#HEAD:<identifier of a lexeme>#, \verb#HEAD:$X# or
	      \verb#HEAD:_# for a predicate, a predicate stored in a variable, or
	      an anonymous predicate.
	\item \verb#LEMMA:<identifier>#, \verb#LEMMA:$X# for a lemma constraint.
	\item \verb#FORM:"<form>"#, \verb#FORM:$X# or \verb#FORM:_# for a literal
	      form, a form stored in a variable, or an anonymous form.  Literal
	      forms may be concatenated with \verb#+#, for example
	      \verb#FORM:"New" + " " + "York"#.
	\item a pair $attr_i:val_i$
	      \begin{itemize}
		      \item with $attr_i$ a named attribute
		      \item and with $val_i$ the value of an attribute that can be:

		            \begin{itemize}
		            \item A named variable
		            \item An anonymous variable
		            \item A nil value (\texttt{NIL})
		            \item An atomic constant
		            \item A set of atomic constants separated by \verb#|#
		            \item A literal constant: integer or double
		            \item A feature structure
		            \item A list
		            \item An inherited or synthesized attribute: $\uparrow$,
		                  $\Downarrow_i$
		            \item A field access such as $\uparrow$\verb#.attr# or
		                  $\Downarrow$\verb#i.attr#
		            \end{itemize}
	      \end{itemize}

\end{itemize}

The lists are noted
\begin{itemize}
	\item $<f_1, f_2, \dots, f_k>$ with $f_i$ as variables, constants,
	      numbers, feature structures, lists, inherited or synthesized
	      attributes, or field accesses.
	\item $<head::tail>$ with $head$ as one list element and $tail$ as a list.
	      $head$ refers to the first item on the list, $tail$ to the rest of
	      the list.  The tail may also be \verb#NIL#.
\end{itemize}

If one wishes to bring an operational semantics to the syntagmatic
rule, one defines $<operational~semantics>$ which is a list of rules
between braces \verb#{...}#.

A rule may optionally be preceded by rule modifiers such as
\verb#@trace#, \verb#@withSpaces#, \verb#@withoutSpaces# or
\verb#@unordered#.

\begin{itemize}
	\item \verb#@trace# traces the application of the rule.
	\item \verb#@withSpaces# generates spaces between the components of the rule.
	\item \verb#@withoutSpaces# suppresses spaces between the components of the rule.
	\item \verb#@unordered# allows the components of the rule to be generated in any order.
\end{itemize}

Each rule is written:

\begin{itemize}

	\item $\{<$list of rules $>\}$ A list of rules.

	\item \verb#attest <expr>;# Explicit guard where \verb#<expr># is an
	      boolean expression that must be checked.

	\item \verb#[f_1, feat_2, \dots, f_k];# Guard where the feature
	      structure must subsume $\uparrow$. The guard may contain the feature
	      \verb#HEAD# which represents the concept to be generated. Thus, a
	      concept is not necessarily linked to a lexical entry with
	      \texttt{Elvex}, unlike usual text generation systems;
	      it can produce a syntactic constraint. Conversely, a syntactic constraint
	      does not necessarily have to be linked to a specific concept.

	\item \verb#print <expr_1>, ..., <expr_n>;# Displays one or several
	      expressions at standard output (useful for debugging).

	\item \verb#println <expr_1>, ..., <expr_n>;# Does the same and displays a
	      new line.

	\item \verb#eprint <expr_1>, ..., <expr_n>;# Displays one or several
	      expressions on standard error (useful for debugging without changing
	      generated output).

	\item \verb#eprintln <expr_1>, ..., <expr_n>;# Does the same on standard
	      error and displays a new line.

	\item  \verb#<expr_1> = <expr_2>;#

	      An assignment of the expression $<expr_1>$ with the value of the
	      expression $<expr_2>$. The expression $<expr_1>$ is a complex
	      variable, or simple variable, the expression $<expr_2>$ is a constant
	      or a complex variable or a simple variable such that there is a
	      environment \footnote{An environment is an application $\phi$ from
		      $E$ to $F$ such that $E$ is a set of variables and $F$ a set of
		      constants or variables. The resolution of a variable $\alpha$ in
		      an environment is $\beta = \phi^k(\alpha)$ such that $\beta$ is a
		      constant. It should be noted that \texttt{Elvex} also handles
		      complex variables. In this case $\phi^k(\alpha)$ is replacing
		      variables in terms by their constants. } of the syntagmatic rule
	      to solve it.

	      The assignments of $\uparrow$ and $\Downarrow_i$ are not
	      possible, as these attributes are implicitly computed during text
	      generation and modifying them would result in
	      inconsistencies.

	      The assignments are limited to these cases:
	      \begin{itemize}

		      \item Assign a list

		            \begin{itemize}
			            \item  \verb#<#$\dots$\verb#> = <#$\dots$\verb#># Assigns a list to
			                  another one.
			            \item \verb#<#$\dots$\verb#> = $X# Assigned by the value of the
			                  variable \verb#$X#
		            \end{itemize}

		      \item Assign an inherited attribute

		            \begin{itemize}
			            \item $\downarrow$\verb#i = $X# Assigned by the value of the
			                  variable \verb#$X#
			            \item $\downarrow$\verb#i = [#$\dots$\verb#]# Assigned by a
			                  constant or variable feature structure.
			            \item $\downarrow$\verb#i = #$\uparrow$ Assigned by an inherited attribute.
			            \item $\downarrow$\verb#i = #$\dots$\verb# #$\cup$\verb# #$\dots$
			                  Assigned by the result of the unification of two expressions
			                  (feature structures or attributes).
			            \item $\downarrow$\verb#i = #$\Downarrow$\verb#j# Assigned by a
			                  synthesized attribute.
			            \item $\downarrow$\verb#i = #$\uparrow$\verb#.attr# or
			                  $\downarrow$\verb#i = #$\Downarrow$\verb#j.attr# assigns the value obtained
			                  by field access.
		            \end{itemize}

		      \item Assign a synthesized attribute

		            \begin{itemize}
			            \item $\Uparrow$\verb# = $X#
			            \item $\Uparrow$\verb# = [#$\dots$\verb#]#
			            \item $\Uparrow$\verb# = #$\uparrow$
			            \item $\Uparrow$\verb# = #$\dots$\verb# #$\cup$\verb# #$\dots$
			            \item $\Uparrow$\verb# = #$\Downarrow$\verb#j#
			            \item $\Uparrow$\verb# = #$\uparrow$\verb#.attr# or
			                  $\Uparrow$\verb# = #$\Downarrow$\verb#j.attr#
		            \end{itemize}

		      \item  Assign a simple variable

		            \begin{itemize}
			            \item \verb#$X = $Y#
			            \item \verb#$X = a# Assigned by a constant or a literal.
			            \item \verb#$X = <#$\dots$\verb#>#
			            \item \verb#$X = [#$\dots$\verb#]#
			            \item \verb#$X = #$\uparrow$
			            \item \verb#$X = #$\dots$\verb# #$\cup$\verb# #$\dots$
			            \item \verb#$X = #$\Downarrow$\verb#j#
			            \item \verb#$X = #$\uparrow$\verb#.attr# or
			                  \verb#$X = #$\Downarrow$\verb#j.attr# by field access
			            \item \verb#$X = <expr>#  by evaluating an arithmetical or logical expression
		            \end{itemize}

		      \item Subsume a feature structure

		            A feature structure is subsumed

		            \begin{itemize}
			            \item \verb#[#$\dots$\verb#] #$\subset$\verb# #$\uparrow$ by an inherited attribute.
			            \item \verb#[#$\dots$\verb#] #$\subset$\verb# #$\Downarrow$\verb#j# by a synthesized attribute
			            \item \verb#[#$\dots$\verb#] #$\subset$\verb# $X# by the value of a variable.
			            \item \verb#[#$\dots$\verb#] #$\subset$\verb# #$\uparrow$\verb#.attr# or
			                  \verb#[#$\dots$\verb#] #$\subset$\verb# #$\Downarrow$\verb#j.attr# by field access.
		            \end{itemize}

		      \item \verb#if (<boolExpr>) <rules>#

		            Interprets operational semantics \verb#<rules># if and only
		            if the expression \verb#<boolExpr># is verified.

		      \item \verb#if (<boolExpr>) <rules_1> else <rules_2>#

		            Interprets operational semantics \verb#<rules_1># if and
		            only if \verb#<boolExpr># is verified, otherwise interpret \verb#<rules_2>#.

		      \item \verb#deferred (<boolExpr>) <rules>#

		            Delays the interpretation of \verb#<rules># until \verb#<boolExpr># can be
		            evaluated. This is useful when a statement depends on synthesized
		            attributes that are not immediately available.

		      \item \verb#order i < j < ...;#

		            Adds a linearization constraint between the components of the
		            right-hand side of the current syntagmatic rule. The indexes are
		            one-based: \verb#1# refers to the first component, \verb#2# to
		            the second one, and so on. For instance, \verb#order 2 < 1;#
		            states that the second component must be generated before the
		            first one. The constraint is projected onto the components that
		            are actually present: if an optional component is absent, the
		            remaining components keep the relative order specified by the
		            chain.  The constraint is a generation and forest-ordering
		            constraint, not a syntactic left-to-right parsing requirement:
		            alternatives that cannot be generated are simply discarded, and
		            the surviving alternatives may still satisfy the order constraint.

		      \item \verb#order << i;#

		            Constrains component \verb#i#, if present, to be generated before
		            all other present components of the rule.

		      \item \verb#order >> i;#

		            Constrains component \verb#i#, if present, to be generated after
		            all other present components of the rule.

		      \item \verb#order by <key_1, key_2, ..., key_k>;#

		            Sorts the present components of the rule according to the keys
		            associated with them.  The first key is associated with the first
		            component, the second key with the second component, and so on.
		            Keys must be numeric and the current parser accepts only a
		            variable, a field access, an integer or a double as a key; it
		            does not accept an arbitrary algebraic expression in this
		            position.  Components are generated in increasing key order;
		            equal keys keep the original right-hand-side order.  For example:
		            \verb#order by <$R1,$R2,$R3>;# or
		            \verb#order by <#$\Downarrow$\verb#1.rank, #$\Downarrow$\verb#2.rank>;#.


	      \end{itemize}

	      A logical expression \verb#<boolExpr># is

	      \begin{itemize}
		      \item \verb#<boolExpr>#  $\lor$ \verb#<boolExpr># Or
		      \item \verb#<boolExpr>#  $\land$ \verb#<boolExpr># And
		      \item $\neg$ \verb#<boolExpr># Negation
		      \item \verb#<boolExpr>#  $\Rightarrow$ \verb#<boolExpr>#
		            Material implication
		      \item \verb#<boolExpr>#  $\Leftrightarrow$ \verb#<boolExpr>#
		            Biconditional
		      \item \verb#<expr>#  $==$ \verb#<expr># Equality
		      \item \verb#<expr>#  $\neq$ \verb#<expr># Difference
		      \item \verb#<expr># Expression evaluated as Boolean: A
		            feature structure is evaluated as true if it is not
		            \verb#NIL# or $\bot$ (result of the failure of unification
		            between two feature structures).
	      \end{itemize}

	      The expressions are:

	      \begin{itemize}
		      \item Integer evaluated as true if not zero
		      \item Double evaluated as true if not zero
		      \item String evaluated as true if not empty
		      \item \verb#<expr>#  $\cup$ \verb#<expr># Unification of two
		            feature structures
		      \item $\uparrow$
		      \item $\Uparrow$
		      \item $\downarrow_i$
		      \item $\Downarrow_i$
		      \item $\uparrow$\verb#.attr# and $\Downarrow$\verb#i.attr# Field access
		      \item \verb|#i| and \verb|#i.j| Presence tests for the generated
		            components identified by one-based indexes
		      \item \verb#[...]# Feature structure
		      \item \verb#NIL# Null value of a feature
		      \item \verb#a|b...# Constant evaluated as true
		      \item \verb#$X# Variable
		      \item \verb#<...># List evaluated as true if not null (\verb#<>#)
	      \end{itemize}

\paragraph{Presence tests: \texttt{\#i} and \texttt{\#i.j}}

The expressions \verb|#i| and \verb|#i.j| test the presence of a component on
the right-hand side of the current rule. Indexes are one-based.

\begin{itemize}
	\item \verb|#i| is true if the $i$-th component is present in the current
	      branch. It is false if this component is an optional component that
	      has been skipped.
	\item \verb|#i.j| is true if the $i$-th component is present and its
	      selected variant is the $j$-th alternative. It is false for the other
	      variants and for an absent optional component.
\end{itemize}

For example, with the rule:

\begin{lstlisting}[numbers=none]
S → (NP|cln) VP { ... }
\end{lstlisting}

\verb|#1| is true when either \verb#NP# or \verb#cln# is present,
\verb|#1.1| is true only when \verb#NP# is selected, and \verb|#1.2| is true
only when \verb#cln# is selected.

Presence tests may be temporarily undecidable while the generator has not yet
resolved the corresponding optional term or alternative. In that case, a
statement depending on the test is suspended and is tried again when the
choice is known. This applies in particular to \verb#if# statements:

\begin{lstlisting}[numbers=none]
S → (A) B {
  if (#1) {
    ↓2 = [x:A];
  }
  else {
    ↓2 = [x:noA];
  }
}
\end{lstlisting}

The \verb#if# is not evaluated on the unresolved item \verb#S → • (A) B#.
It is evaluated only after the optional component has been resolved: the
present branch makes \verb|#1| true, and the absent branch makes it false.
For a mandatory simple component, \verb|#i| is true. For a mandatory
alternative, \verb|#i| is true once one variant has been selected, and
\verb|#i.j| identifies the selected variant.

\paragraph{Predefined string constants}

The lexer provides two predefined string constants that are useful for
debugging:

\begin{itemize}
	\item \verb#__LINE__# expands to the current line number;
	\item \verb#__FILE__# expands to the current file name.
\end{itemize}

They are treated as strings and may be used in expressions, for instance
with printing statements:

\begin{lstlisting}[numbers=none]
eprintln("debug at " + __FILE__ + ":" + __LINE__);
\end{lstlisting}

Let's add algebraic expressions and numbers:

	      \begin{itemize}
		      \item \verb#<expr>#  $<$ \verb#<expr>#
		      \item \verb#<expr>#  $\le$ \verb#<expr>#
		      \item \verb#<expr>#  $>$ \verb#<expr>#
		      \item \verb#<expr>#  $\ge$ \verb#<expr>#
		      \item \verb#<expr>#  $+$ \verb#<expr>#
		      \item \verb#<expr>#  $-$ \verb#<expr>#
		      \item \verb#<expr>#  $*$ \verb#<expr>#
		      \item \verb#<expr>#  \% \verb#<expr>#
		      \item \verb#<expr>#  $/$ \verb#<expr>#
		      \item $-$ \verb#<expr>#
		      \item \verb#floor(<expr>)# Integer floor of a numeric expression
		      \item \verb#rand()# Random number generator
		      \item Double-precision floating-point format number.  In the current
		            lexer, double literals are written with a leading \verb#d#,
		            for example \verb#d1.25#.
	      \end{itemize}

\end{itemize}

\subsubsection*{Variable evaluation}

Variables are evaluated locally at a syntagmatic rule. This
defines the scope of a variable.

Here are different methods for solving variables:

\begin{itemize}
	\item Direct assignment
	      \begin{lstlisting}
S → S {
  [vform:NIL];
  [tense:$tense, aspect:$aspect, mode:$mode] ⊂ ↑;
  if (!$tense)
    $tense = present;
...
}
\end{lstlisting}
	      Line 5 indicates that the value of the variable \verb#$tense# is
	      \verb#present# by default (in the absence of another value previously given). The
	      value of
	      \verb#$tense# is eventually given by the inherited attribute
	      $\uparrow$ on line 3.

	      Note that the order of evaluation of the rules is indifferent and that it
	      is not necessary to precede rule (3) in the order that they are listed.

	\item Guard
	      \begin{lstlisting}[numbers=none]
S → NP VP {
  [subject:$inheritedSubject, $Rest];
...
}
\end{lstlisting}

	      The value of the variable \verb#$inheritedSubject# is given
	      by the $\uparrow$ subject feature. The value of the variable
	      \verb#$Rest# is the feature structure that remains when you remove the
	      feature \verb#subject#.

	\item Subsumption
	      \begin{lstlisting}[numbers=none]
S → NP VP {
...
  [subject:$synthesizedSubject] ⊂ ⇓2;
...
}
\end{lstlisting}

	      The value of the variable \verb#$synthesizedSubject# is given
	      by the \verb#subject# feature of $\Downarrow2$.

	\item Binding a variable within a syntagmatic rule
	      \begin{lstlisting}[numbers=none]
S → NP VP {
  [subject:$inheritedSubject, $Rest];
  [subject:$synthesizedSubject] ⊂ ⇓2;
  ↓1 = [subject:$inheritedSubject] 
  		∪ [subject:$synthesizedSubject];
...
}
\end{lstlisting}

	      The variables \verb#$synthesizedSubject# and \verb#$inheritedSubject#
	      are both used within this rule to build the \verb#subject# feature of
	      the inherited attribute $\downarrow1$.

	\item Linking a variable within a lexical entry

	      \begin{lstlisting}[numbers=none]
a	v[HEAD:brouillard, subjC:[FORM:"il"], 
		objC:[HEAD:brouillard, det:yes, 
		def:no, part:yes, neg:$Neg], 
		locCl:[HEAD:_pro, number:sg, 
			person:three], 
		vform:tensed, vtense:present, 
		mode:indicative, 
		subj:[person:three, number:sg], 
		neg:$Neg, fct:none]
\end{lstlisting}

	      The variable \verb#$Neg# is linked in this lexical entry to the expression "it is foggy"
	      so that its negative use propagates a negation to the noun:

	      \begin{exe}
		      \ex \begin{xlist}
			      \ex \label{0-a} \textit{Il y a du brouillard} (it is foggy)
			      \ex \label{0-a} \textit{Il n'y a pas de brouillard} (There is no fog)
			      \ex \label{0-a} \textit{*Il y a de brouillard} ()
			      \ex \label{0-a} \textit{*Il n'y a pas du brouillard}
		      \end{xlist}
	      \end{exe}

\end{itemize}


\subsection{Macros}

Macro files contain definitions of the form:

\begin{lstlisting}[numbers=none]
@name = [feature:value, ...];
\end{lstlisting}

A macro expands to the feature structure on the right-hand side.  It can be
used as an item in a feature structure, either alone or mixed with explicit
features:

\begin{lstlisting}[numbers=none]
@commonNoun = [cat:noun, number:sg];

N → n {
  ↓1 = [@commonNoun, HEAD:$H];
}
\end{lstlisting}

If a macro name is unknown, parsing fails with an error.  Macro expansion is
syntactic: the features stored in the macro are inserted where \verb#@name#
appears.

\subsection{Lexicon}

The lexicon is a list of entries separated by semicolons.

Each entry defines the sequence to be produced, followed by a disjunction
(noted by the character \verb#|#) of lexical categories
optionally followed by a feature structure.

The sequence to be produced may be:

\begin{itemize}
	\item A string of characters between \verb#""#, in which case the sequence
	      will be exactly this one. Note that the sequence may be empty.

	\item A string of characters between \verb#""# containing
	      variables. The sequence produced will be this one with the variables
	      substituted by their respective values.

	\item \texttt{FORM}: the sequence produced will literally be the one
	      given by the feature \texttt{FORM}; variables in the form are resolved
	      during generation.

	\item \texttt{FORM} followed by an identifier in a lexical entry, for
	      example \verb#FORM title;#, declares an entry whose produced form is
	      supplied dynamically.

\end{itemize}

The category is an identifier, namely the final symbol of grammar.

The feature structure allows to represent all the
constraints imposed by the lexicon.

A few examples:

\begin{itemize}
	\item \verb#"?" interrogativeDot;#

	      \bigskip
	\item \verb#"tag-$Line" tag[line:$Line]; #

	      \bigskip
	\item \verb# FORM title; #

	      \bigskip
	\item \verb#"elles-mêmes"	itself [gender:fm, number:pl, person:three, itself:yes]; #

	      \bigskip
	\item \begin{verbatim}
serai	aux_être [aux:être, voice:active, finite:yes,
  mode:indicative,  vtense:future_anterieur, vform:tensed, 
subj:[person:one, number:sg]]; 
\end{verbatim}

	      \bigskip
	\item \begin{verbatim}
// Lexical input "temperature" producing 
//"La température est de xxx degrés" (the temperature is xxx degrees Celsius"
est v[HEAD:température, 
	subjC:[HEAD:température, number:sg, det:yes, def:yes], 
	pobjC:[HEAD:degré, number:pl, pcas:de, det:yes,
		detNum:[HEAD:num, value:$Deg]], 
	vform:tensed, vtense:present, mode:indicative, 
	subj:[person:three, number:sg], fct:none, value:$Deg]

//"la température est douce" (the temperature is mild)
|v[HEAD:température, subjC:[HEAD:température, number:sg, 
	det:yes, def:yes], modC:[HEAD:doux, number:sg, gender:fm],
	vform:tensed, vtense:present, mode:indicative, 
	subj:[person:three, number:sg], fct:few]

//"la température est élevée" (the temperature is high)
|v[HEAD:température, subjC:[HEAD:température, number:sg, 
	det:yes, def:yes], modC:[HEAD:élevé, number:sg, gender:fm],
	vform:tensed, vtense:present, mode:indicative, 
	subj:[person:three, number:sg], fct:much];
\end{verbatim}

	      \bigskip
	\item \begin{verbatim}
// Lexical entry of "temperature".
// producing "Il fait chaud" (It's hot)
fait	v[HEAD:température, subjC:[FORM:"il"], 
	modC:[HEAD:chaud, number:sg, gender:ms], vform:tensed,
	vtense:present, mode:indicative, 
	subj:[person:three, number:sg], neg:$Neg, fct:high]

//"Il fait très chaud" (It's very hot)
|v[HEAD:température, subjC:[FORM:"il"], modC:[HEAD:chaud, 
	number:sg, gender:ms, mod:<[HEAD:très]>], vform:tensed, 
	vtense:present, mode:indicative, subj:[person:three, 
	number:sg], neg:$Neg, fct:very_high]

//"Il fait très froid" (It's very cold)
|v[HEAD:température, subjC:[FORM:"il"], modC:[HEAD:froid, 
	number:sg, gender:ms, mod:<[HEAD:très]>], vform:tensed,
	vtense:present, mode:indicative, subj:[person:three, 
	number:sg], neg:$Neg, fct:very_low]

//"Il fait froid" (It's cold)
|v[HEAD:température, subjC:[FORM:"il"], modC:[HEAD:froid, 
	number:sg, gender:ms], vform:tensed, vtense:present, 
	mode:indicative, subj:[person:three, number:sg], 
	neg:$Neg, fct:low]; 
\end{verbatim}

	      \bigskip
	\item \begin{verbatim}
// Lexical entry of "rain" producing
// "Il y a une pluie fine" (A fine rain fell)
a	v[HEAD:pleuvoir, subjC:[FORM:"il"], 
	objC:[HEAD:pluie, det:yes, number:sg, qual:yes,
		def:no, mod:<[HEAD:fin, pos:post]>], 
	locCl:[HEAD:_pro, number:sg, person:three], 
	vform:tensed, vtense:present, mode:indicative, 
	subj:[person:three, number:sg], neg:$Neg, fct:few]
	
// "il y a un peu de pluie" (It is raining a bit)
|v[HEAD:pleuvoir, subjC:[FORM:"il"], 
	objC:[HEAD:pluie, det:yes, number:sg, 
	detForm:[FORM:"un peu de"]], 
	locCl:[HEAD:_pro, number:sg, person:three], 
	vform:tensed, vtense:present, mode:indicative, 
	subj:[person:three, number:sg], neg:$Neg, fct:not_much]

// "il y a beaucoup de pluie" (There is plenty of rainfall)
|v[HEAD:pleuvoir, subjC:[FORM:"il"], 
	objC:[HEAD:pluie, det:yes, number:sg, 
	detForm:[FORM:"beaucoup de"]], 
	locCl:[HEAD:_pro, number:sg, person:three], 
	vform:tensed, vtense:present, mode:indicative, 
	subj:[person:three, number:sg], neg:$Neg, fct:much]

// "il y a de fortes précipitations" (It is raining so much)
|v[HEAD:pleuvoir, subjC:[FORM:"il"], 
	objC:[HEAD:précipitation,
        det:yes, number:pl, qual:yes, def:no, mod:<[HEAD:fort,
        pos:pre]>], locCl:[HEAD:_pro, number:sg, person:three],
        vform:tensed, vtense:present, mode:indicative,
        subj:[person:three, number:sg], neg:$Neg, fct:a_lot]; 
\end{verbatim}

\end{itemize}

\subsection{Input}

The input is summarized by the type of phrase or text you want to generate, followed by a
feature structure containing the concept and the illocutionary structure to be expressed.


\texttt{HEAD} allows you to define the named concept to be generated.
The special features \texttt{HEAD}, \texttt{LEMMA} and \texttt{FORM} can also
be accessed by field access, for example $\uparrow$\verb#.HEAD#, $\uparrow$\verb#.LEMMA# or
$\Downarrow$\verb#1.FORM#.  This is useful when a rule has to inspect the concept,
lemma or literal form stored in an inherited or synthesized feature
structure.

examples

\begin{itemize}
	\item
	      \begin{verbatim}
sentence [HEAD:raining, neg:yes]
\end{verbatim}

	      Simple entry for the concept of \textit{raining} producing
	      \begin{exe}
		      \ex \begin{xlist}
			      \ex \textit{Il ne pleut pas} (It is not raining)
			      \ex \textit{Il n'y a pas de pluie} (There is no rain)
		      \end{xlist}
	      \end{exe}


	\item
	      \begin{verbatim}
sentence 
   [HEAD:to_ask,
    i:[HEAD:menuisier, id:3, number:sg, gender:ms, def:yes],
    iii:[HEAD:apprenti, id:6, def:yes, gender:ms, number:sg],
    ii:[HEAD:scier,
        i:[idref:6],
        ii:[HEAD:poutre, number:sg, def:yes]
        ]
    ]
\end{verbatim}

	      Complex concept where the features \texttt{id} and \texttt{idref}
	      make it possible to constrain co-references.

	      This entry produces
	      \begin{exe}
		      \ex
		      \begin{tabbing}
			      \textit{Le menuisier} \= \textit{demande} \= \textit{à} \textit{l'apprenti}~~~ \= \textit{de} \textit{scier} \= \textit{la} \textit{poutre}\\
			      The carpenter \> asks \> the apprentice \> to saw \> the beam\\
			      The carpenter asks the apprentice to saw the
			      beam
		      \end{tabbing}
	      \end{exe}

	\item
	      \begin{verbatim}
sentence 
   [HEAD:raining, 
    tense:future, 
    modSType:time, 
    modS:<[FORM:"lundi 18 mars 2019", type:time]>
   ]
\end{verbatim}

	      Concept producing an adverbial on the basis of the
	      constant "lundi 18 mars 2019" (Monday, the 18th of March 2019). This entry produces:

	      \begin{exe}
		      \ex \textit{Il va pleuvoir le lundi 18 mars 2019.} (It will rain on Monday, March 18, 2019)
	      \end{exe}

	\item
	      \begin{verbatim}
sentence 
   [HEAD:cause, 
    i: [HEAD:raining]]
    ii: [HEAD:se_couvrir, i:[HEAD:you], modality:devoir]
   ]
\end{verbatim}

	      Input for the following generated sentences:
	      \begin{exe}
		      \ex \begin{xlist}
			      \ex \textit{Il pleut, tu dois te couvrir} (It is raining, you have to dress warmly)
			      \ex \textit{Il pleut, par conséquent tu dois te couvrir} (It is raining, therefore you have to dress warmly)
			      \ex \textit{Tu dois te couvrir s'il pleut} (You have to dress warmly if it is raining)
			      \ex \textit{Tu dois te couvrir car il pleut} (You have to dress warmly because it is raining)
			      \ex \textit{S'il pleuvait, tu devrais te couvrir} (If it rained, you should dress warmly)
			      \ex \textit{S'il pleut, tu devras te couvrir} (If it is raining, you should dress warmly)
			      \ex \textit{Il pleuvrait, tu devrais te couvrir} (It would rain, you should dress warmly)
		      \end{xlist}
	      \end{exe}

	\item
	      \begin{verbatim}
sentence 
   [HEAD:initiative, 
    i:[HEAD:Jean], 
    objC:[lexFct:Magn, mod:<[HEAD:audacieux]>], 
    lexFct:Oper1
   ]
\end{verbatim}

	      Input containing features to fix lexical functions. The
	      production is
	      \begin{exe}
		      \ex \textit{Jean prend une initiative très audacieuse.} (Jean takes a very bold initiative)
	      \end{exe}

	\item
	      \begin{verbatim}
sentence 
   [HEAD:initiative, 
    i:[HEAD:Jean], 
    objC:[lexFct:Magn], 
    lexFct:Oper1
   ]
\end{verbatim}

	      Input for

	      \begin{exe}
		      \ex \textit{Jean prend une belle initiative.}  (Jean takes a fine initiative)
	      \end{exe}

	\item
	      \begin{verbatim}
sentence 
   [HEAD:to_break,
    i:[HEAD:Jean], 
    ii:[HEAD:carafe, number:sg, def:yes, 
       mod:<[HEAD:little],[HEAD:white]>], 
    diathesis:passive, tense:past, neg:yes, 
    modV:<[HEAD:roughly]>
   ]
\end{verbatim}

	      Input containing a negative passive feature. The
	      production is
	      \begin{exe}
		      \ex \textit{La petite carafe blanche n'avait pas été cassée
			      brutalement par Jean.} (The small white carafe had not been broken
		      brutally by Jean.)
	      \end{exe}

\end{itemize}

\section{Design pattern (to be completed)}

A design pattern applies in a general case which can be summarized by
some constraints. It is not associated with a specific event, but to a
recurring organization of the grammar for which a single and coherent
answer must be provided.

\subsubsection*{A - Local propagation of synthesized attributes}

Context:
A feature \texttt{A} that must constrain a result \texttt{R}, where
\texttt{R} depends only on the lexicon or grammar.

Examples of applications:
\begin{itemize}
	\item Production of a constrained noun in an idiomatic phrase
	      (\textit{breaking the \textbf{ice}})
	\item Conjugation to the passive voice
\end{itemize}

The general rule is written as follows

\begin{lstlisting}
X → Y {
  [A:$v, $rest];
  [ARealization:R] ⊂ ⇓1;
  ↓1 = [R:$v, $rest];
}
\end{lstlisting}

Example 1: The sentence \textit{We basically got taken to the cleaners
	on that deal.} introduces the fixed noun \textit{cleaners} which
comes from the verb \textit{take one to the cleaner}. Here the agent of the
passive sentence is not represented.

\begin{lstlisting}
VN → V (PP) PP {
  [obj:NIL, obl:$Obl];
  [subj:$subj] ⊂ ↑;
  ↓1 = ↑;
  ↓3 = $Obl;
  [oblC:$oblSynt] ⊂ ⇓1;
  if (#2) {
    attest $oblSynt;
    ↓2 = $oblSynt;
  }
  else {
    attest ¬$oblSynt;
  }
}

taken verb[HEAD:take_one_to_the_cleaners, 
	constI:[HEAD:cleaner, number:pl, def:yes]];
\end{lstlisting}

Example 2:
The agent of a clause is carried out by the grammatical function subject
for conjugation to the active voice, and by the oblique function in
\textit{by} for a passive voice conjugation.

\begin{lstlisting}
// agent as subject
Sentence → Sentence {
  [agent:$agent, $rest];
  [agentRealization:subject] ⊂ ⇓1;
  ↓1 = [subject:$agent, $rest];
  ⇑ = ⇓1;
}

// agent as oblique
Sentence → Sentence {
  [agent:$agent, $rest];
  [agentRealization:oblBy] ⊂ ⇓1;
  ↓1 = [byObj:$agent, $rest];
  ⇑ = ⇓1;
}

// patient as object
Sentence → Sentence {
  [patient:$patient, agent:NIL, $rest];
  [patientRealization:object] ⊂ ⇓1;
  ↓1 = [object:$patient, $rest];
  ⇑ = ⇓1;
}

// patient as subject
Sentence → Sentence {
  [patient:$patient, agent:NIL, $rest];
  [patientRealization:subject] ⊂ ⇓1;
  ↓1 = [subject:$patient, $rest];
  ⇑ = ⇓1;
}

// passive voice
Sentence → NP VERB PP {
  [subject:$subject, parObl:$parObl, $rest];
  ↓1 = $subject;
  ↓2 = $rest ∪ [voice:passive];
  ↓3 = $parObl;
  ⇑ = [agentRealization:oblBy, 
	patientRealization:subject];
}

// active voice
Sentence → NP VERB NP {
  [subject:$subject, object:$object, $rest];
  ↓1 = $subject;
  ↓2 = $rest ∪ [voice:active];
  ↓3 = $object;
  ⇑ = [agentRealization:subject, 
	patientRealization:object];
}
\end{lstlisting}

\subsubsection*{B - Transformation of the attribute inherited by a syntactic head}

Context:
An inherited attribute \texttt{A} that must constrain a realization
\texttt{R} by a new predicate in generation.

Examples of applications:
\begin{itemize}
	\item Use of a verbal tense auxiliary or a modal auxiliary
	\item Compound conjugation in French
\end{itemize}

The general rule is written as follows

\begin{lstlisting}
X → Y {
  [A:K];
  [HEAD:$HEAD, A:k, $Rest] ⊂ ↑;
  ↓1 = [HEAD:L, arg:[HEAD:$HEAD, $Rest]];
}

\end{lstlisting}

Example:
In French, the aspect is reflected in conjugated forms, but also in
compound verb constructions (\textit{être en train de} + infinitive, to be in the
process of doing sth).

\begin{lstlisting}
sentence → sentence {
  [aspect:progressive];
  [HEAD:$HEAD, aspect:progressive, i:$I, $Rest] ⊂ ↑;
  ↓1 = [HEAD:être_en_train_de, 
	i:[HEAD:$HEAD, i:$I, $Rest]];
  ⇑ = ⇓1;
}
\end{lstlisting}

\section{The \texttt{elvex} command}

\texttt{Elvex} is a Natural Language Generator. It takes as input a
concept, a local lexicon, a compacted lexicon, and a grammar, and then
outputs a text corresponding to the concept written in natural language.

\subsection*{Compilation and full installation of \texttt{Elvex}}

The full compilation of \texttt{Elvex} requires the following programs on your machine:

\begin{itemize}
	\item g++
	\item flex
	\item bison
	\item cmake
	\item libxml2-dev
\end{itemize}

The full command is:

\begin{lstlisting}[frame=none, numbers=none]
cmake .
make
sudo make install
\end{lstlisting}

\subsection{\texttt{Elvex} command}

The \texttt{elvex} command is:

\begin{lstlisting}[frame=none, numbers=none]
elvex [options] [<input>]*
\end{lstlisting}

The options are:

\begin{description}
	\item \verb#-h|--help#

	      Displays help and exits

	\item \verb#-v|--version#

	      Displays the version number and exits

	\item \verb#-V|--verbose#

	      Enables verbose mode.

	\item \verb#-a|--reduce-all#

	      Reduces all rules during the generation process. This mode
	      allows fine debugging and requires an understanding of the
	      generation process to be mastered.

	\item \verb#-t|--trace#

	      Displays the reduced rules during the generation process. This mode is useful
	      for grammar debugging.

	\item \verb#-r|--random#

	      Displays only one randomly selected output. Random choice may occur at
	      several levels: lexical entry selection, forest node selection and final
	      root forest selection.

	\item \verb#--seed <number>#

	      Sets the random seed for reproducible random generation. This option is
	      useful with \verb#--random# and also controls calls to \verb#rand()# in
	      operational statements.

	\item \verb#-f|--first#

	      Displays only the first output in the current deterministic generation
	      order. Unlike \verb#--random#, this option does not shuffle choices.

	\item \verb#--strategy <exhaustive|sample|beam>#

	      Selects the derivation strategy. The default strategy is \verb#exhaustive#.

	\item \verb#--max-rule-choices <integer>#

	      Defines the maximum number of rule or disjunction choices opened in
	      \verb#sample# mode.

	\item \verb#--beam-width <integer>#

	      Defines the maximum number of items kept per state in \verb#beam# mode.

	\item \verb#--macros-file <file>#

	      File containing macro definitions.  Each definition has the form
		      \verb#@name = [features];#.

	\item \verb#--rules-file <file>#

	      File containing the rules.

	\item \verb#--lexicon-file <file>#

	      File containing the local lexicon.

	\item \verb#--input-file <file>#

	      File containing the input.

	\item \verb#--compacted-lexicon-file <path/file># or \verb#--clf <path/file>#

	      Prefix of the compacted lexicon files. The files
	      \verb#<path/file>.tbl# and \verb#<path/file>.fsa# are used.

	\item \verb#--xml <file>#

	      Writes an XML representation of the generation forest to
	      \verb#<file>#. This option is available only when \texttt{Elvex}
	      has been compiled with XML output support.

	\item \verb#--max-length <integer>#

	      Defines the maximum length of the sequence to be produced.

	\item \verb#--max-usages <integer>#

	      Defines the maximum number of times the same rule can be used.

	\item \verb#--max-items <integer>#

	      Defines the maximum cardinality of the item sets.

	\item \verb#--max-time <integer>#

	      Defines the number of seconds before the program stops.

	\item \verb#--max-attempts <integer>#

	      Defines the maximum number of attempts for random generation.

\end{description}

When a grammar can generate the empty string, the corresponding realization
contains no surface form. Depending on the output context, this may be shown
as an empty line or may be visually invisible among the other realizations.


\subsection{Grammar debugging}

Grammar debugging is performed with the \texttt{elvex} command using
trace options. For example:

\begin{lstlisting}[frame=none, numbers=none]
elvex --trace \
  --rules-file rules.elvex \
  --lexicon-file lexicon.elvex \
  <input>
\end{lstlisting}

With tracing enabled, \texttt{elvex} produces an HTML document on the
standard output. You must use a browser to view this document properly
and review the computation.

Where the specific options are:

\begin{description}

	\item \verb#--trace-all#

	      Trace all generation steps

	\item \verb#--trace-init#

	      Trace the initial items.

	\item \verb#--trace-stage#

	      Trace the stabilized item set for each generation stage. A stage is
	      printed once after normalization and closure; this is different from
	      \verb#--trace-close#, which traces the internal closure operations.

	\item \verb#--trace-close#

	      Trace the item set closure.

	\item \verb#--trace-shift#

	      Trace the items for each lexical entry.

	\item \verb#--trace-reduce#

	      Trace the items for each reduction.

	\item \verb#--trace-action#

	      Trace the operational semantics results.

\end{description}


\subsection{\texttt{elvexlexicon} for managing compacted lexicons}

The \texttt{elvexlexicon} command builds and inspects compacted lexicons.
Compacted lexicons are intended for open lexical categories such as nouns,
verbs, adjectives and adverbs. Closed or grammatical categories, such as
prepositions, determiners, pronouns or complementizers, may remain in the
local lexicon loaded directly by \texttt{elvex} with \verb#--lexicon-file#.

The command is:

\begin{lstlisting}[frame=none, numbers=none]
elvexlexicon [global options] <build|consult|list>
\end{lstlisting}

The action is mandatory:

\begin{description}
	\item \verb#build#

	      Builds the compacted lexicon from a pattern file and a morphological
	      file.

	\item \verb#list#

	      Lists the content of an existing compacted lexicon.

	\item \verb#consult#

	      Consults an existing compacted lexicon interactively.
\end{description}

The global options are:

\begin{description}
	\item \verb#-h|--help#

	      Displays help and exits.

	\item \verb#-v|--version#

	      Displays the version number and exits.

	\item \verb#-V|--verbose#

	      Enables verbose mode.

	\item \verb#--compacted-lexicon-file <path/file># or \verb#--clf <path/file>#

	      Prefix of the compacted lexicon files. The files
	      \verb#<path/file>.tbl# and \verb#<path/file>.fsa# are used.

	\item \verb#--macros-file <file>#

	      File containing macro definitions used in the pattern and
		      morphological files.  Each definition has the form
		      \verb#@name = [features];#.

	\item \verb#--pattern-file <file>#

	      Pattern file describing lexical entries at the lexeme level. It
	      consists of lines of the form:

	      \begin{lstlisting}[frame=none, numbers=none]
lexeme  pos  lemma  features
\end{lstlisting}

	      The \verb#lexeme# is the lexical unit selected by the generation
	      model. It may correspond to a simple word, a collocation or an
	      idiomatic construction. The \verb#lemma# identifies the morphological
	      paradigm with which the lexeme is associated.

	\item \verb#--morpho-file <file>#

	      Morphological file describing inflected forms at the lemma level. It
	      consists of lines of the form:

	      \begin{lstlisting}[frame=none, numbers=none]
form  pos  lemma  features
\end{lstlisting}

	      The \verb#form# is the produced surface form. The \verb#lemma# is used
	      to match a morphological entry with one or more pattern entries.
\end{description}

The compacted lexicon is built by combining pattern entries and morphological
entries that share the same \verb#pos# and \verb#lemma#. Their feature
structures are unified. This keeps a clear distinction between the lexeme,
which represents a lexical choice in the generation model, and the lemma,
which represents a morphological paradigm.

For example, a pattern file may contain adjective classes and their
linearization rank:

\begin{lstlisting}[frame=none, numbers=none]
LARGE  adj  large  [HEAD:large, type:size, rank:1]
PINK   adj  pink   [HEAD:pink, type:color, rank:4]
IRON   adj  iron   [HEAD:iron, type:material, rank:6]
\end{lstlisting}

and the morphological file may contain the corresponding surface forms:

\begin{lstlisting}[frame=none, numbers=none]
large  adj  large  []
pink   adj  pink   []
iron   adj  iron   []
\end{lstlisting}

After compilation, the resulting entries can be used by grammar rules. For
instance, the synthesized feature \verb#rank# can later be exploited by a
linearization instruction such as:

\begin{lstlisting}[frame=none, numbers=none]
order by <⇓1.rank, ⇓2.rank, ⇓3.rank>;
\end{lstlisting}

\subsubsection*{Command to build a compacted lexicon}

\begin{lstlisting}[frame=none, numbers=none]
elvexlexicon --clf <path/file> \
  --macros-file <macros-file> \
  --pattern-file <pattern-file> \
  --morpho-file <morpho-file> \
  build
\end{lstlisting}

The \verb#--macros-file# option is needed only if the pattern or
morphological files use macros.

\subsubsection*{Command to list a compacted lexicon}

\begin{lstlisting}[frame=none, numbers=none]
elvexlexicon --clf <path/file> list
\end{lstlisting}

\subsubsection*{Command to consult a compacted lexicon}

\begin{lstlisting}[frame=none, numbers=none]
elvexlexicon --clf <path/file> consult
\end{lstlisting}

The compacted lexicon is a set of lexical entries of the type

\paragraph{}\verb@pos#lexeme	form#FS|form#FS...@

\begin{itemize}
	\item \verb#pos# is a part of speech
	\item \verb#lexeme# is the lexeme given by the \verb#HEAD# feature
	\item \verb#form# is the produced form, represented as a string of UTF-8 characters
	\item \verb#FS# is a feature structure
\end{itemize}

\end{document}
